%% ch06.tex --- Logistic Regression
%% Numbers are computed by verify/ch06_logistic.py.
%%
%% Note on ordering: BOOK_PLAN lists "why squared error fails" first, but the
%% non-convexity argument needs the sigmoid model to exist before it can be
%% made.  The chapter therefore splits that item in two --- the failure of
%% linear regression (Sec. 6.2, before the sigmoid) and the failure of squared
%% error on the sigmoid model (Sec. 6.5, after it).

\chapter{Logistic Regression}
\label{ch:ch06}

\epigraph{It is the mark of a truly intelligent person to be moved by
statistics.}{---\ George Bernard Shaw}

\begin{chapterlearning}
After working through this chapter you will be able to:
\begin{enumerate}
  \item Say why a linear model fitted by least squares is the wrong tool for a
        binary target, and demonstrate the failure on data.
  \item State the sigmoid function, derive its derivative, and list the four
        properties that make it the right squashing function.
  \item Read the output of a logistic model as a probability, and convert
        between probability, odds and log-odds.
  \item Find the decision boundary of a logistic model, linear or polynomial,
        and explain why it is a property of the model rather than of the data.
  \item Show that squared error on a sigmoid model gives a non-convex cost, and
        say what that costs an optimizer in practice.
  \item Derive the logistic loss from maximum likelihood, in full, and reduce
        its two cases to one line.
  \item Prove that the logistic cost is convex, and say what the proof requires
        of the data.
  \item Derive the gradient, run the update by hand, and explain why it has the
        same \emph{form} as the linear-regression update without being the same
        algorithm.
  \item Choose a decision threshold from the costs of the two kinds of error
        rather than accepting $0.5$.
\end{enumerate}
\end{chapterlearning}

\Cref{ch:ch01} divided supervised learning into regression and classification
and then spent four chapters on regression. This chapter takes up the other
half. The target is now a category --- malignant or benign, spam or not spam,
fraudulent or legitimate --- and almost everything has to be rebuilt: the model,
the cost, and the way the answer is read. The optimizer, remarkably, does not
change at all.

We build it in the order of \cref{sec:four-parts}: data, hypothesis space, cost,
optimizer. The data comes first, and it comes with a constraint the previous
chapters did not have.

\section{The Classification Problem}
\label{sec:classification-setup}

\begin{defbox}{Binary classification}
A \term{binary classification} problem is a supervised learning problem in which
$\Yspace = \{0,1\}$. The class labelled $1$ is called the \term{positive class}
and the class labelled $0$ the \term{negative class}; the names are conventional
and carry no value judgement --- in medical screening the positive class is
usually the disease.

A learner for such a problem is a \term{classifier}. It must ultimately emit a
label in $\{0,1\}$, but as this chapter shows, the useful thing to learn first is
a number in $[0,1]$.
\end{defbox}

The labels $0$ and $1$ are arbitrary tags. Nothing in the problem says that
malignant is one more than benign, and the pitfall in \cref{sec:supervised}
warned against reading arithmetic into such codes. The whole difficulty of this
chapter follows from that: we are about to build a model that outputs a real
number, and we must be careful about what the number means.

Throughout the chapter we use a small screening dataset.

\begin{defbox}{Dataset T7}
Eight tumours, each recorded by its greatest diameter in centimetres, each
labelled $1$ if malignant and $0$ if benign.

\begin{center}
\small
\renewcommand{\arraystretch}{1.15}
\begin{tabular}{@{}l*{8}{r}@{}}
\toprule
$x$ (cm) & $1$ & $2$ & $3$ & $4$ & $5$ & $6$ & $7$ & $8$\\
$y$      & $0$ & $1$ & $0$ & $0$ & $0$ & $1$ & $1$ & $1$\\
\bottomrule
\end{tabular}
\end{center}

Larger tumours are more often malignant, but the rule is not perfect: the $2$~cm
tumour was malignant and the $5$~cm one was not. That imperfection is
deliberate, and \cref{sec:gd-logistic} shows what a dataset without it would do
to the fit.
\end{defbox}

\section{Why Linear Regression Fails}
\label{sec:linreg-fails}

The obvious first attempt is to ignore the difficulty. The target is a number,
$0$ or $1$; fit a straight line by least squares as in \cref{ch:ch02}, and
predict malignant whenever the line is above $0.5$. On T7 this gives
\[
  f(x) = 0.119048\,x - 0.035714 ,
\]
which crosses $0.5$ at $x = 4.5$ and gets six of the eight labels right. It looks
like a working classifier. It is not, for two reasons.

\subsection{The output is not a probability, or anything else}

The fitted values on the eight training points run from $0.0833$ to $0.9167$,
which is at least in range. Step outside the data and the model reports
\[
  f(0) = -0.0357, \qquad f(12) = 1.3929, \qquad f(25) = 2.9405 .
\]
A tumour of zero size is $-3.6\%$ malignant and one of $25$~cm is $294\%$
malignant. There is no reading of the output under which those are answers. The
model has been asked for a quantity confined to $[0,1]$ and given a hypothesis
space --- straight lines --- in which no member is confined to $[0,1]$.

\subsection{The fit is driven by distance, not by correctness}

The second failure is worse, because it corrupts the part of the model that did
work. Least squares minimises squared \emph{distance} to the targets. A training
example far out along the $x$-axis, correctly classified and never in doubt,
still has a large residual if the line does not pass close to it, and the fit
bends to reduce that residual.

Add to T7 a single further malignant tumour, of size $20$~cm, whose class no
method could get wrong, and refit --- with least squares and, for comparison,
with the logistic model this chapter is about to build:

\begin{center}
\small
\renewcommand{\arraystretch}{1.2}
\begin{tabular}{@{}lrrrr@{}}
\toprule
 & \multicolumn{3}{c}{least squares} & logistic\\
\cmidrule(lr){2-4}\cmidrule(lr){5-5}
extra example & $b$ & $w$ & boundary & boundary\\
\midrule
none      & $-0.035714$ & $0.119048$ & $4.5000$ & $4.5000$\\
$x = 10$  & $ 0.003226$ & $0.108065$ & $4.5970$ & $4.4644$\\
$x = 15$  & $ 0.164286$ & $0.069048$ & $4.8621$ & $4.4978$\\
$x = 20$  & $ 0.266087$ & $0.046522$ & $5.0280$ & $4.4999$\\
$x = 30$  & $ 0.362366$ & $0.026344$ & $5.2245$ & $4.5000$\\
$x = 50$  & $ 0.427509$ & $0.013400$ & $5.4097$ & $4.5000$\\
$x = 100$ & $ 0.467576$ & $0.005822$ & $5.5691$ & $4.5000$\\
\bottomrule
\end{tabular}
\end{center}

The least-squares boundary walks from $4.50$ to $5.57$ as the extra tumour
recedes, driven entirely by a point about which there was never any question.
The logistic boundary does not move: it settles at $4.5000$ and stays there. The
same effect flattens the line --- with the extra point at $x = 20$ the eight
original fitted values compress into $[0.3126,\,0.6383]$, so the model is now
mildly unsure about everything, including the $1$~cm tumour it previously called
benign with confidence.

\begin{figure}[htbp]
  \centering
  \begin{minipage}[t]{0.48\linewidth}
    \centering
    \begin{tikzpicture}
      \begin{axis}[mlaxis, width=\linewidth, height=4.6cm,
        xlabel={tumour size $x$ (cm)}, ylabel={$f(x)$},
        xmin=0, xmax=9.4, ymin=-0.15, ymax=1.15,
        legend style={font=\tiny, at={(0.98,0.06)}, anchor=south east}]
        \addplot[MLAmber, thick, dashed, domain=0:9.4] {0.119048*x - 0.035714};
        \addlegendentry{least squares}
        \addplot[MLBlue, very thick, domain=0:9.4, samples=120]
          {1/(1+exp(-(0.594084*x - 2.673380)))};
        \addlegendentry{logistic}
        \addplot[MLGray, dotted, thick, forget plot, domain=0:9.4] {0.5};
        \addplot[only marks, mark=*, mark size=1.7pt, MLRose, forget plot]
          coordinates {(1,0)(3,0)(4,0)(5,0)};
        \addplot[only marks, mark=square*, mark size=1.7pt, MLTeal, forget plot]
          coordinates {(2,1)(6,1)(7,1)(8,1)};
      \end{axis}
    \end{tikzpicture}
    {\scriptsize Both cross $0.5$ at $x = 4.5$.\par}
  \end{minipage}\hfill
  \begin{minipage}[t]{0.48\linewidth}
    \centering
    \begin{tikzpicture}
      \begin{axis}[mlaxis, width=\linewidth, height=4.6cm,
        xlabel={tumour size $x$ (cm)}, ylabel={$f(x)$},
        xmin=0, xmax=21.5, ymin=-0.15, ymax=1.15,
        legend style={font=\tiny, at={(0.98,0.06)}, anchor=south east}]
        \addplot[MLAmber, thick, dashed, domain=0:21.5]
          {0.046522*x + 0.266087};
        \addlegendentry{least squares}
        \addplot[MLBlue, very thick, domain=0:21.5, samples=160]
          {1/(1+exp(-(0.594373*x - 2.674608)))};
        \addlegendentry{logistic}
        \addplot[MLGray, dotted, thick, forget plot, domain=0:21.5] {0.5};
        \addplot[only marks, mark=*, mark size=1.7pt, MLRose, forget plot]
          coordinates {(1,0)(3,0)(4,0)(5,0)};
        \addplot[only marks, mark=square*, mark size=1.7pt, MLTeal, forget plot]
          coordinates {(2,1)(6,1)(7,1)(8,1)(20,1)};
      \end{axis}
    \end{tikzpicture}
    {\scriptsize One extra example at $x = 20$ moves the line, not the curve.\par}
  \end{minipage}
  \caption{Least squares against logistic regression on T7. Circles are benign,
  squares malignant. The straight line leaves $[0,1]$ on both sides, and a
  single distant example that no method would misclassify drags its threshold
  crossing from $4.50$ to $5.03$. The logistic curve is bounded, and its
  crossing stays at $4.4999$, because an example it already predicts correctly
  and confidently contributes almost nothing to the gradient.}
  \label{fig:linreg-fails}
\end{figure}

\begin{keyinsight}{The two requirements}
A classifier needs a hypothesis space whose members map $\R$ into $[0,1]$, and a
cost that charges for being \emph{wrong} rather than for being \emph{far}. Linear
regression fails both tests. The rest of this chapter supplies one replacement
for each: the sigmoid for the first, the logistic loss for the second.
\end{keyinsight}

\section{The Sigmoid Function}
\label{sec:sigmoid}

Keep the linear part $z = \vw\cdot\vx + b$, which carries all the parameters, and
pass it through a function that squashes $\R$ onto $(0,1)$.

\begin{defbox}{The sigmoid, and the logistic model}
The \term{sigmoid} (or \term{logistic}) function is
\[
  \gsig(z) \;=\; \frac{1}{1+e^{-z}}, \qquad \gsig:\R \to (0,1).
\]
The \term{logistic regression} model composes it with a linear function:
\[
  \fvwb(\vx) \;=\; \gsig(\vw\cdot\vx + b) \;=\; \frac{1}{1+e^{-(\vw\cdot\vx+b)}} .
\]
The parameters are the same $\vw \in \Real{n}$ and $b \in \R$ as in
\cref{ch:ch04}. Only the last step is new.
\end{defbox}

Four properties do the work, and all four are worth verifying rather than
memorising.

\begin{enumerate}
  \item \textbf{Range.} $e^{-z} > 0$ for every real $z$, so the denominator
        exceeds $1$ and $0 < \gsig(z) < 1$. The endpoints are approached but
        never attained: a logistic model is never completely certain.
  \item \textbf{Centre.} $\gsig(0) = 1/2$. The point $z = 0$ is where the model
        has no opinion, and it will turn out to be the decision boundary.
  \item \textbf{Symmetry.} $1 - \gsig(z) = \gsig(-z)$. Reversing the sign of
        every parameter swaps the two classes, so the labelling of positive and
        negative is genuinely arbitrary.
  \item \textbf{Monotonicity.} $\gsig$ is strictly increasing, so ordering by
        $\gsig(z)$ is the same as ordering by $z$. This is what makes the
        decision boundary of \cref{sec:decision-boundary} linear.
\end{enumerate}

\begin{derivation}{The derivative of the sigmoid}
Write $\gsig(z) = (1+e^{-z})^{-1}$ and differentiate by the chain rule:
\[
  \gsig'(z) \;=\; -(1+e^{-z})^{-2}\cdot\bigl(-e^{-z}\bigr)
            \;=\; \frac{e^{-z}}{(1+e^{-z})^{2}} .
\]
Split the fraction into two factors, one of which is $\gsig(z)$:
\[
  \gsig'(z) \;=\; \underbrace{\frac{1}{1+e^{-z}}}_{\gsig(z)}
                 \cdot\underbrace{\frac{e^{-z}}{1+e^{-z}}}_{?} .
\]
The second factor simplifies by adding and subtracting $1$ in its numerator:
\[
  \frac{e^{-z}}{1+e^{-z}} = \frac{(1+e^{-z}) - 1}{1+e^{-z}}
    = 1 - \frac{1}{1+e^{-z}} = 1 - \gsig(z).
\]
Hence
\[
  \boxed{\;\gsig'(z) = \gsig(z)\bigl(1-\gsig(z)\bigr)\;}
\]
The derivative is expressible in the function's own value, which is why
implementations cache the forward pass and why the same identity reappears,
unchanged, in the backpropagation of \cref{ch:ch09}.

Two consequences follow immediately. Since $\gsig \in (0,1)$, the product
$\gsig(1-\gsig)$ is strictly positive, so $\gsig$ is strictly increasing. And
$p(1-p)$ is maximised at $p = 1/2$, so
\[
  \max_z \gsig'(z) = \tfrac12\cdot\tfrac12 = \tfrac14 \quad\text{at } z=0,
\]
falling away to zero in both tails. \Cref{sec:squared-error-fails} shows what
that vanishing derivative does to a badly chosen cost.
\end{derivation}

\begin{figure}[htbp]
  \centering
  \begin{tikzpicture}
    \begin{axis}[mlaxis, width=0.74\linewidth, height=4.8cm,
      xlabel={$z$}, ylabel={},
      xmin=-6.4, xmax=6.4, ymin=-0.06, ymax=1.12,
      legend style={font=\scriptsize, at={(0.02,0.98)}, anchor=north west}]
      \addplot[MLBlue, very thick, domain=-6.4:6.4, samples=160]
        {1/(1+exp(-x))};
      \addlegendentry{$\gsig(z)$}
      \addplot[MLAmber, thick, dashed, domain=-6.4:6.4, samples=160]
        {1/(1+exp(-x)) * (1 - 1/(1+exp(-x)))};
      \addlegendentry{$\gsig'(z) = \gsig(1-\gsig)$}
      \addplot[MLGray, dotted, thick, forget plot, domain=-6.4:6.4] {0.5};
    \end{axis}
  \end{tikzpicture}
  \caption{The sigmoid and its derivative. The curve passes through
  $(0,\tfrac12)$, saturates towards $0$ and $1$, and its slope peaks at
  $\tfrac14$ exactly where it crosses $\tfrac12$. Beyond $|z| = 4$ the slope is
  under $0.018$: the model is not merely confident there, it is nearly
  unresponsive.}
  \label{fig:sigmoid}
\end{figure}

\begin{center}
\footnotesize
\setlength{\tabcolsep}{4pt}
\renewcommand{\arraystretch}{1.15}
\begin{tabular}{@{}l*{7}{r}@{}}
\toprule
$z$        & $-6$ & $-4$ & $-2$ & $0$ & $2$ & $4$ & $6$\\
\midrule
$\gsig(z)$  & $0.002473$ & $0.017986$ & $0.119203$ & $0.500000$
            & $0.880797$ & $0.982014$ & $0.997527$\\
$\gsig'(z)$ & $0.002467$ & $0.017663$ & $0.104994$ & $0.250000$
            & $0.104994$ & $0.017663$ & $0.002467$\\
\bottomrule
\end{tabular}
\end{center}

\subsection{The output as a probability}

The sigmoid guarantees a number in $(0,1)$, which permits a probabilistic
reading, but does not by itself justify one. The justification is the modelling
assumption we are about to make, and it is worth stating explicitly because
everything in \cref{sec:logistic-loss} is derived from it.

\begin{defbox}{The logistic model's assumption}
\[
  \Prob(y = 1 \mid \vx;\ \vw,b) \;=\; \fvwb(\vx),
  \qquad
  \Prob(y = 0 \mid \vx;\ \vw,b) \;=\; 1 - \fvwb(\vx).
\]
The two exhaust the possibilities and sum to $1$, so one number determines both.
Read $\fvwb(\vx) = 0.7$ as ``a $70\%$ chance that this example is positive'',
\emph{given that the assumption holds}.
\end{defbox}

The assumption is an assumption, not a theorem, and a fitted model can be
confident and wrong. \Cref{ch:ch10} shows how to check whether the numbers a
model calls probabilities behave like probabilities.

\subsection{Odds and log-odds}

The reason the sigmoid appears here rather than some other bounded increasing
curve is best seen by inverting it. Write $p = \gsig(z)$ and solve for $z$:
\[
  p(1+e^{-z}) = 1 \;\Longrightarrow\;
  e^{-z} = \frac{1-p}{p} \;\Longrightarrow\;
  z = \ln\!\left(\frac{p}{1-p}\right).
\]
The quantity $p/(1-p)$ is the \term{odds} of the positive class, and its
logarithm is the \term{log-odds}, or \term{logit}. So the model says
\[
  \ln\frac{\Prob(y=1\mid\vx)}{\Prob(y=0\mid\vx)} \;=\; \vw\cdot\vx + b .
\]

\begin{keyinsight}{Logistic regression is linear regression on the log-odds}
The model is not a mangled version of linear regression. It is an exactly linear
model of a different quantity: not $y$, and not $\Prob(y=1)$, but the logarithm
of the odds. Everything that made linear models tractable --- one parameter per
feature, additive effects, a convex fitting problem --- survives on that scale.

The reading of a coefficient follows. Adding one to $x_j$ adds $w_j$ to the
log-odds, hence multiplies the odds by $e^{w_j}$. The fitted model of
\cref{sec:gd-logistic} has $w = 0.594084$ per centimetre, so each additional
centimetre of tumour multiplies the odds of malignancy by
$e^{0.594084} = 1.8114$ --- a little under a doubling. Note that the effect on
the \emph{probability} is not constant: it is largest near $p = 1/2$ and
vanishes in the tails, which is the sigmoid's shape restated.
\end{keyinsight}

\section{Decision Boundaries}
\label{sec:decision-boundary}

A probability is not yet a decision. To classify, compare it with a threshold;
take $0.5$ for now and revisit the choice in \cref{sec:threshold}.

Because $\gsig$ is increasing and $\gsig(0) = \tfrac12$,
\[
  \fvwb(\vx) \ge \tfrac12
  \;\Longleftrightarrow\;
  \vw\cdot\vx + b \ge 0 ,
\]
so the sigmoid disappears from the decision entirely.

\begin{defbox}{Decision boundary}
The \term{decision boundary} of a logistic model at threshold $\tfrac12$ is the
set
\[
  \{\vx : \vw\cdot\vx + b = 0\},
\]
a hyperplane in $\Real{n}$: a point when $n=1$, a line when $n=2$, a plane when
$n=3$. It is a property of the \emph{parameters}, not of the training data ---
the same boundary can be drawn before any data is seen, by choosing $\vw$ and
$b$ by hand.
\end{defbox}

With one feature the boundary is the single point $x^\star = -b/w$. For T7 the
fitted parameters are $b = -2.673380$ and $w = 0.594084$, so
$x^\star = 2.673380/0.594084 = 4.5000$ cm: predict malignant above $4.5$~cm,
benign below.

\begin{workedexample}{A boundary in two dimensions}
Dataset~T8 records twenty tumours by size $x_1$ (cm) and a texture score $x_2$
on a $0$--$10$ scale.

\begin{center}
\footnotesize
\renewcommand{\arraystretch}{1.1}
\begin{tabular}{@{}lrrrrrrrrrr@{}}
\toprule
\multicolumn{11}{@{}l}{\textbf{Benign} ($y=0$)}\\
$x_1$ & $2.4$ & $1.8$ & $2.2$ & $2.9$ & $3.8$ & $2.0$ & $2.9$ & $2.4$ & $3.0$ & $1.9$\\
$x_2$ & $3.4$ & $0.4$ & $4.6$ & $2.7$ & $5.9$ & $3.2$ & $3.1$ & $5.0$ & $2.7$ & $2.7$\\
\midrule
\multicolumn{11}{@{}l}{\textbf{Malignant} ($y=1$)}\\
$x_1$ & $3.8$ & $3.5$ & $3.8$ & $2.4$ & $3.3$ & $3.2$ & $3.9$ & $3.9$ & $4.4$ & $4.9$\\
$x_2$ & $6.3$ & $7.2$ & $6.2$ & $4.4$ & $6.7$ & $8.2$ & $6.5$ & $6.8$ & $5.5$ & $5.2$\\
\bottomrule
\end{tabular}
\end{center}

Fitting gives $b = -9.827806$, $w_1 = 0.931932$, $w_2 = 1.362882$, with cost
$J = 0.262254$ and eighteen of twenty labels correct. The boundary is the line
$0.931932\,x_1 + 1.362882\,x_2 - 9.827806 = 0$, which is easier to plot solved
for $x_2$:
\[
  x_2 \;=\; 7.2110 - 0.6838\,x_1 .
\]
At $x_1 = 1.5$ it passes through $x_2 = 6.1854$; at $x_1 = 5.0$, through
$x_2 = 3.7921$. The two mistakes are exactly the two points on the wrong side of
that line: the benign tumour at $(3.8, 5.9)$, which the model calls malignant
with probability $0.8525$, and the malignant one at $(2.4, 4.4)$, which it calls
benign with probability $1 - 0.1688 = 0.8312$. Both are confidently wrong, and
no straight line can rescue either.
\end{workedexample}

\begin{figure}[htbp]
  \centering
  \begin{minipage}[t]{0.48\linewidth}
    \centering
    \begin{tikzpicture}
      \begin{axis}[mlaxis, width=4.5cm, height=4.5cm, scale only axis,
        xlabel={size $x_1$}, ylabel={texture $x_2$},
        xmin=1.4, xmax=5.3, ymin=0, ymax=9]
        \addplot[MLBlue, very thick, domain=1.4:5.3, forget plot]
          {7.2110 - 0.6838*x};
        \addplot[only marks, mark=*, mark size=1.6pt, MLRose, forget plot]
          coordinates {(2.4,3.4)(1.8,0.4)(2.2,4.6)(2.9,2.7)(3.8,5.9)
                       (2.0,3.2)(2.9,3.1)(2.4,5.0)(3.0,2.7)(1.9,2.7)};
        \addplot[only marks, mark=square*, mark size=1.6pt, MLTeal, forget plot]
          coordinates {(3.8,6.3)(3.5,7.2)(3.8,6.2)(2.4,4.4)(3.3,6.7)
                       (3.2,8.2)(3.9,6.5)(3.9,6.8)(4.4,5.5)(4.9,5.2)};
      \end{axis}
    \end{tikzpicture}
    {\scriptsize Linear: $\vw\cdot\vx + b = 0$.\par}
  \end{minipage}\hfill
  \begin{minipage}[t]{0.48\linewidth}
    \centering
    \begin{tikzpicture}
      \begin{axis}[mlaxis, width=4.5cm, height=4.5cm, scale only axis,
        xlabel={$x_1$}, ylabel={$x_2$},
        xmin=-2.2, xmax=2.2, ymin=-2.2, ymax=2.2]
        \addplot[MLBlue, very thick, domain=0:360, samples=120, forget plot]
          ({cos(x)}, {sin(x)});
        \addplot[only marks, mark=*, mark size=1.6pt, MLRose, forget plot]
          coordinates {(0,0)(0.5,0.3)(-0.4,0.5)(0.3,-0.6)(-0.6,-0.3)(0.15,0.7)};
        \addplot[only marks, mark=square*, mark size=1.6pt, MLTeal, forget plot]
          coordinates {(1.6,0.4)(-1.5,0.9)(0.9,-1.5)(-0.8,-1.6)(1.2,1.3)
                       (-1.7,-0.5)(0.2,1.8)(1.5,-0.7)};
      \end{axis}
    \end{tikzpicture}
    {\scriptsize Quadratic: $x_1^2 + x_2^2 - 1 = 0$.\par}
  \end{minipage}
  \caption{Two decision boundaries. Left: T8 with its fitted line
  $x_2 = 7.2110 - 0.6838\,x_1$; the two misclassified points are the two on the
  wrong side. Right: the boundary produced by $\vw = (0,0,1,1)$ and $b=-1$ on
  the features $(x_1,x_2,x_1^2,x_2^2)$ --- a circle, from a model that is still
  linear in its parameters.}
  \label{fig:boundary}
\end{figure}

\subsection{Non-linear boundaries}

A hyperplane in the space of \emph{features} need not be a hyperplane in the
space of the original measurements, and \cref{sec:feature-engineering} already
supplied the mechanism. Give the model the four features
$(x_1, x_2, x_1^2, x_2^2)$ and choose
\[
  \vw = (0,\,0,\,1,\,1), \qquad b = -1 .
\]
Then $z = x_1^2 + x_2^2 - 1$, and the boundary $z = 0$ is the unit circle:
inside it the model predicts $0$, outside it predicts $1$. Higher-degree terms
buy ellipses, then curves with several disconnected pieces, then essentially any
shape at all.

The parameters above were chosen, not fitted, so that the boundary is exactly
the unit circle; fitting them would be the identical procedure on the identical
machinery, since the model remains linear in $\vw$ and $b$.

\begin{pitfall}{More flexibility, same trap}
The picture of \cref{ch:ch05} repeats here. Enough polynomial features will
separate any training set perfectly, boundary contorted to enclose each stray
point, and the training cost will be excellent. \Cref{ch:ch07} shows what
happens on new data, and adds the penalty term that is the standard defence ---
for logistic regression exactly as for linear regression.
\end{pitfall}

\section{Why Squared Error Fails}
\label{sec:squared-error-fails}

The hypothesis space is now correct: every member of it maps into $(0,1)$. It is
tempting to keep the cost of \cref{ch:ch02} and write
\[
  J_{\mathrm{sq}}(\vw,b)
  \;=\; \frac{1}{2m}\sum_{i=1}^{m}\bigl(\gsig(\vw\cdot\xii{i}+b) - \yii{i}\bigr)^2 .
\]
This is where the second requirement of \cref{sec:linreg-fails} bites. The cost
is bounded, well defined, and differentiable --- and it is a bad objective, for
two independent reasons.

\subsection{It is not convex}

\Cref{sec:convexity} proved that the squared-error cost of linear regression is
convex, and \cref{ch:ch03} leaned on that fact throughout: convexity is why any
local minimum is global, and why gradient descent cannot be trapped. Composing
the model with a sigmoid destroys it. Convexity requires a positive
semi-definite Hessian everywhere; on T7,
\begin{center}
\small
\renewcommand{\arraystretch}{1.2}
\begin{tabular}{@{}lrrr@{}}
\toprule
$(b,w)$ & $J_{\mathrm{sq}}$ & $\lambda_{\min}$ & $\lambda_{\max}$\\
\midrule
$(0.00,\ 0.00)$  & $0.125000$ & $+0.012476$ & $+1.643774$\\
$(2.00,\ -0.80)$ & $0.237035$ & $-0.134306$ & $+0.016137$\\
$(4.00,\ -1.00)$ & $0.276621$ & $-0.331318$ & $-0.000688$\\
\bottomrule
\end{tabular}
\end{center}
One negative eigenvalue anywhere is enough to settle the question, and the third
row has two: at $(b,w) = (4,-1)$ the surface is concave in \emph{every}
direction, which is something no convex function does anywhere. Every guarantee
\cref{ch:ch03} attached to convexity is therefore withdrawn --- a local minimum
need no longer be global, and where a run ends may depend on where it began.

\subsection{Its gradient vanishes where you most need it}

The second reason is the more practical, and it follows from
$\max \gsig' = \tfrac14$. Differentiating the per-example squared error with
respect to $z$ gives, by the chain rule,
\[
  \frac{\partial}{\partial z}\,\tfrac12\bigl(\gsig(z)-y\bigr)^2
  \;=\; (f - y)\,f(1-f), \qquad f = \gsig(z).
\]
The factor $f(1-f)$ is the sigmoid's derivative, and it collapses to nothing in
both tails. A model that is confidently, catastrophically wrong sits in a tail:

\begin{center}
\small
\renewcommand{\arraystretch}{1.2}
\begin{tabular}{@{}lrrrr@{}}
\toprule
$f$ (with $y=1$) & squared loss & log loss
                 & $\partial(\text{sq})/\partial z$
                 & $\partial(\text{log})/\partial z$\\
\midrule
$0.9$   & $0.005000$ & $0.105361$ & $-0.009000$ & $-0.100000$\\
$0.5$   & $0.125000$ & $0.693147$ & $-0.125000$ & $-0.500000$\\
$0.1$   & $0.405000$ & $2.302585$ & $-0.081000$ & $-0.900000$\\
$0.01$  & $0.490050$ & $4.605170$ & $-0.009801$ & $-0.990000$\\
$0.001$ & $0.499001$ & $6.907755$ & $-0.000998$ & $-0.999000$\\
\bottomrule
\end{tabular}
\end{center}

Read down the squared-error column. A prediction of $0.001$ when the truth is
$1$ --- as wrong as a probability can be --- is penalised $0.499$, barely more
than the $0.405$ charged for a mild error of $0.1$, and \emph{less} than four
times what a coin flip costs. Worse, its gradient is $0.000998$: the algorithm
receives almost no signal to fix it. The log-loss column, defined in the next
section, charges $6.91$ and delivers a gradient of $0.999$, a thousand times
larger.

\begin{figure}[htbp]
  \centering
  \begin{minipage}[t]{0.48\linewidth}
    \centering
    \begin{tikzpicture}
      \begin{axis}[mlaxis, width=\linewidth, height=4.4cm,
        xlabel={$w$}, ylabel={$J_{\mathrm{sq}}$},
        xmin=-2.6, xmax=3.1, ymin=0.05, ymax=0.28]
        \addplot[MLRose, very thick] coordinates {
          (-2.50,0.2499)(-2.25,0.2499)(-2.00,0.2498)(-1.75,0.2497)
          (-1.50,0.2496)(-1.25,0.2493)(-1.00,0.2489)(-0.75,0.2480)
          (-0.50,0.2462)(-0.25,0.2409)(0.00,0.2198)(0.25,0.1516)
          (0.50,0.0881)(0.75,0.1047)(1.00,0.1411)(1.25,0.1666)
          (1.50,0.1806)(1.75,0.1879)(2.00,0.1924)(2.25,0.1962)
          (2.50,0.2000)(2.75,0.2040)(3.00,0.2084)};
      \end{axis}
    \end{tikzpicture}
    {\scriptsize Squared error: a $0.0003$ drop over $w \in [-2.5,-1.5]$.\par}
  \end{minipage}\hfill
  \begin{minipage}[t]{0.48\linewidth}
    \centering
    \begin{tikzpicture}
      \begin{axis}[mlaxis, width=\linewidth, height=4.4cm,
        xlabel={$w$}, ylabel={$J_{\mathrm{log}}$},
        xmin=-2.6, xmax=3.1, ymin=0, ymax=9]
        \addplot[MLBlue, very thick] coordinates {
          (-2.50,8.5250)(-2.25,7.8065)(-2.00,7.0880)(-1.75,6.3697)
          (-1.50,5.6517)(-1.25,4.9339)(-1.00,4.2167)(-0.75,3.5006)
          (-0.50,2.7871)(-0.25,2.0813)(0.00,1.4034)(0.25,0.8346)
          (0.50,0.5483)(0.75,0.5725)(1.00,0.7703)(1.25,1.0505)
          (1.50,1.3734)(1.75,1.7208)(2.00,2.0835)(2.25,2.4565)
          (2.50,2.8365)(2.75,3.2217)(3.00,3.6108)};
      \end{axis}
    \end{tikzpicture}
    {\scriptsize Log loss: a useful slope everywhere.\par}
  \end{minipage}
  \caption{The same slice through both costs on T7, at $b = -2.673380$. Note the
  vertical scales: the squared-error curve spans $0.20$ while the log-loss curve
  spans $8.0$. Left of $w = -1$ the squared-error cost is flat to three decimal
  places, and the gradient there has magnitude $0.0026$ falling to $0.00013$ ---
  a run started in that region makes no measurable progress for thousands of
  iterations.}
  \label{fig:nonconvex}
\end{figure}

Both defects show up in a single experiment. Start both costs at
$(b,w) = (4,-1)$ --- a model that believes larger tumours are safer --- and run
gradient descent with $\alpha = 0.2$ on each:

\begin{center}
\footnotesize
\renewcommand{\arraystretch}{1.2}
\begin{tabular}{@{}rrlrl@{}}
\toprule
iteration & $J_{\mathrm{sq}}$ & $(b,w)$ & $J_{\mathrm{log}}$ & $(b,w)$\\
\midrule
$1$      & $0.276342$ & $(+3.9950,\,-1.0055)$ & $1.443313$ & $(+4.0121,\,-0.6611)$\\
$10$     & $0.272617$ & $(+3.9488,\,-1.0716)$ & $1.325978$ & $(+3.6057,\,-0.5151)$\\
$100$    & $0.249834$ & $(+3.6886,\,-1.5762)$ & $0.703735$ & $(+0.3489,\,+0.0355)$\\
$1{,}000$  & $0.245891$ & $(+2.9135,\,-1.3909)$ & $0.528102$ & $(-2.6585,\,+0.5913)$\\
$10{,}000$ & $0.082297$ & $(-4.8727,\,+0.9244)$ & $0.528099$ & $(-2.6734,\,+0.5941)$\\
$20{,}000$ & $0.075234$ & $(-8.5693,\,+1.5770)$ & $0.528099$ & $(-2.6734,\,+0.5941)$\\
\bottomrule
\end{tabular}
\end{center}

The log-loss run has essentially converged by iteration $1{,}000$ and is exact
by $10{,}000$. The squared-error run still has the sign of $w$ wrong after a
thousand iterations --- it is crawling across the plateau of
\cref{fig:nonconvex} --- and after $200{,}000$ iterations it has reached
$(b,w) = (-25.02,\,4.56)$ with a gradient norm of $1.8\times 10^{-4}$, still not
converged, its coefficients still growing.

\begin{pitfall}{``Squared error is the cost function''}
It is not. Squared error is one cost, appropriate to one situation: a real-valued
target whose errors are plausibly symmetric and unbounded. The cost is a
modelling choice, made after the target and the hypothesis space are known, and
choosing it badly can leave a perfectly good model unfittable. \Cref{ch:ch09}
makes the same choice again for multi-class outputs, and \cref{ch:ch15} again for
ratings.
\end{pitfall}

\section{The Logistic Loss from Maximum Likelihood}
\label{sec:logistic-loss}

Rather than patch the squared-error cost, derive the right one. The model has
already committed to a probabilistic claim in \cref{sec:sigmoid}; the principle
of \term{maximum likelihood} turns that claim into a cost mechanically: choose
the parameters under which the data actually observed was most probable.

\begin{derivation}{The logistic cost, from the model's own assumption}
\textbf{Step 1: one example, two cases.} The assumption is
\[
  \Prob(y=1\mid\vx) = f, \qquad \Prob(y=0\mid\vx) = 1-f,
  \qquad f = \fvwb(\vx).
\]

\textbf{Step 2: one example, one line.} Since $y$ is $0$ or $1$, the two cases
collapse into a single expression:
\[
  \Prob(y\mid\vx) \;=\; f^{\,y}\,(1-f)^{\,1-y}.
\]
Check it: at $y=1$ this is $f^1(1-f)^0 = f$; at $y=0$ it is $f^0(1-f)^1 = 1-f$.
The exponents act as switches, each turning off the factor that does not apply.

\textbf{Step 3: the whole training set.} Treating the $m$ examples as
independent given their inputs, the probability of observing all the labels we
observed is the product
\[
  \mathcal{L}(\vw,b) \;=\; \prod_{i=1}^{m}
    \bigl(f^{(i)}\bigr)^{\yii{i}}\bigl(1-f^{(i)}\bigr)^{1-\yii{i}},
  \qquad f^{(i)} = \fvwb(\xii{i}).
\]
This is the \term{likelihood}: a function of the parameters, not of the data.

\textbf{Step 4: take logarithms.} The product of $m$ numbers below $1$
underflows for any realistic $m$, and products are awkward to differentiate.
Since $\ln$ is increasing, maximising $\ln\mathcal{L}$ maximises $\mathcal{L}$:
\[
  \ln\mathcal{L}(\vw,b) \;=\; \sum_{i=1}^{m}
    \Bigl[\yii{i}\ln f^{(i)} + (1-\yii{i})\ln\bigl(1-f^{(i)}\bigr)\Bigr].
\]

\textbf{Step 5: negate and average.} Optimizers minimise, and dividing by $m$
makes the cost comparable across training sets of different sizes, exactly as in
\cref{ch:ch02}:
\[
  J(\vw,b) \;=\; -\frac{1}{m}\ln\mathcal{L}(\vw,b)
  \;=\; -\frac{1}{m}\sum_{i=1}^{m}
    \Bigl[\yii{i}\ln f^{(i)} + (1-\yii{i})\ln\bigl(1-f^{(i)}\bigr)\Bigr].
\]
Nothing was chosen along the way except the model's own probabilistic
assumption. The cost was not designed; it was deduced.
\end{derivation}

\begin{defbox}{Logistic loss}
The \term{logistic loss}, or \term{binary cross-entropy}, of a single prediction
is
\[
  \loss(f, y) \;=\;
  \begin{cases}
    -\ln f       & \text{if } y = 1,\\[2pt]
    -\ln(1-f)    & \text{if } y = 0,
  \end{cases}
\]
or, equivalently and in one line,
\[
  \loss(f,y) = -y\ln f - (1-y)\ln(1-f),
\]
and the cost is its average over the training set,
$J(\vw,b) = \frac1m\sum_{i} \loss\bigl(f^{(i)}, \yii{i}\bigr)$.
\end{defbox}

The two-case form is the one to reason with and the one-line form the one to
implement. Read the two cases:

\begin{itemize}
  \item If $y = 1$, the loss is $-\ln f$. It is $0$ when $f = 1$, about $0.69$
        when $f = \tfrac12$, and rises without bound as $f \to 0$: predicting
        $0.001$ for a positive example costs $6.91$, and predicting $0$ exactly
        would cost infinity.
  \item If $y = 0$, the loss is $-\ln(1-f)$, the mirror image.
\end{itemize}

\begin{keyinsight}{The loss is unbounded, and that is the point}
Squared error caps the penalty for a confidently wrong prediction at $\tfrac12$;
the logistic loss lets it grow without limit. That is precisely what makes the
gradient survive into the tails, and it is why a logistic model, unlike the
least-squares line of \cref{fig:linreg-fails}, refuses to be confidently wrong
about a training example if any parameters can avoid it.

A useful reference value: at the origin every $f^{(i)} = \tfrac12$, so
$J = -\ln\tfrac12 = \ln 2 = 0.693147$, whatever the data. A trained classifier
whose cost is not clearly below $0.693$ has learned nothing that a coin could
not have told you.
\end{keyinsight}

\begin{researchspotlight}{Cox (1958) --- The Regression Analysis of Binary Sequences}
\textbf{The problem.} Statisticians had regression for continuous responses and
contingency tables for categorical ones, but no satisfactory way to relate a
binary outcome to several continuous explanatory variables at once. Fitting a
straight line to a $0/1$ response gave fitted values outside $[0,1]$ and
non-constant variance, as \cref{sec:linreg-fails} shows.

\textbf{The idea.} Model not the probability but its \emph{logit},
$\ln\{p/(1-p)\}$, as a linear function of the explanatory variables. The logit
transformation was already in use in bioassay; what was missing, and what Cox
supplied, was its treatment as a regression with several explanatory variables.
The logit maps $(0,1)$ onto the whole real line, so a linear model on that scale
can never produce an impossible probability, and the coefficients read as
multiplicative effects on the odds.

\textbf{The result.} Cox developed the estimation theory for this model --- the
maximum-likelihood derivation of \cref{sec:logistic-loss} --- together with its
large-sample properties and worked applications. Logistic regression became, and
remains, the standard tool for binary outcomes in epidemiology, credit scoring
and the social sciences.

\textbf{Why it is here.} The model of this chapter is Cox's, unchanged. What
machine learning added was scale and an optimizer: where Cox solved the
likelihood equations by Newton's method on a handful of variables, we minimise
the same objective by gradient descent on millions. The single logistic unit is
also the building block of \cref{ch:ch08} --- a neural network is, at its
simplest reading, a stack of these.
\parencite{cox1958}
\end{researchspotlight}

\section{Convexity}
\label{sec:logistic-convexity}

Deducing the cost would be of limited comfort if the result were as hard to
minimise as the one it replaced. It is not.

\begin{derivation}{The logistic cost is convex}
Write $z = \vw\cdot\vx + b$ and $f = \gsig(z)$, and consider the loss on one
example as a function of $z$ alone.

\textbf{Step 1: the first derivative, both cases at once.} For $y = 1$,
$\loss = -\ln \gsig(z)$, so using $\gsig' = \gsig(1-\gsig)$,
\[
  \frac{\dd \loss}{\dd z} = -\frac{\gsig'(z)}{\gsig(z)}
  = -\bigl(1-\gsig(z)\bigr) = f - 1 .
\]
For $y = 0$, $\loss = -\ln\bigl(1-\gsig(z)\bigr)$, so
\[
  \frac{\dd \loss}{\dd z} = \frac{\gsig'(z)}{1-\gsig(z)} = \gsig(z) = f - 0 .
\]
Both cases are $f - y$. This cancellation --- the sigmoid's derivative in the
numerator against the logarithm's in the denominator --- is the reason the two
choices belong together, and it is why the gradient of
\cref{sec:gd-logistic} is so simple.

\textbf{Step 2: the second derivative.}
\[
  \frac{\dd^2 \loss}{\dd z^2} = \frac{\dd f}{\dd z}
  = f(1-f) \;>\; 0 \quad\text{for every } z,
\]
strictly, since $0 < f < 1$. So $\loss$ is a strictly convex function of $z$.

\textbf{Step 3: from $z$ to the parameters.} The map $(\vw,b)\mapsto z$ is
affine, and a convex function composed with an affine map is convex. Writing
$\tilde{\vx} = (1, \vx)$ and $\theta = (b,\vw)$ so that $z = \theta\cdot\tilde{\vx}$,
the Hessian of the cost is
\[
  H \;=\; \frac1m\sum_{i=1}^{m} f^{(i)}\bigl(1-f^{(i)}\bigr)\,
          \tilde{\vx}^{(i)}\bigl(\tilde{\vx}^{(i)}\bigr)^{\mathsf{T}} .
\]
For any direction $\mathbf{v}$,
\[
  \mathbf{v}^{\mathsf{T}} H \mathbf{v}
  = \frac1m\sum_{i=1}^{m} f^{(i)}\bigl(1-f^{(i)}\bigr)
    \bigl(\tilde{\vx}^{(i)}\!\cdot\mathbf{v}\bigr)^2 \;\ge\; 0 ,
\]
a sum of non-negative terms. So $H$ is positive semi-definite everywhere and $J$
is convex. It is \emph{strictly} convex --- and the minimum therefore unique ---
exactly when the vectors $\tilde{\vx}^{(i)}$ span $\Real{n+1}$, which is the
same full-rank condition that made $\Xmat^{\mathsf{T}}\Xmat$ invertible in
\cref{sec:singular}.
\end{derivation}

Evaluated over a $201\times201$ grid of $(b,w)$ on T7, the smallest eigenvalue of
the logistic Hessian is $1.3\times10^{-10}$ --- zero to within rounding, and
never negative, in contrast to the $-0.331$ found for squared error.

\begin{keyinsight}{What is preserved, and what is lost}
Convexity means everything \cref{ch:ch03} promised still holds: no local minima
that are not global, no dependence on initialisation, and a divergence threshold
$\alpha < 2/\lambda_{\max}$ that means what it meant before.

One thing is lost. The linear-regression cost was a \emph{quadratic}, which is
why setting its gradient to zero gave the normal equation of
\cref{sec:normal-equation} and a closed-form answer. The logistic cost is convex
but not quadratic: $\grad J = \mathbf{0}$ is a system of transcendental
equations with no solution in elementary functions. For logistic regression
there is no closed form, and iteration is not a convenience but the only option.
\end{keyinsight}

\section{Gradient Descent for Logistic Regression}
\label{sec:gd-logistic}

\begin{derivation}{The gradient}
By Step 1 of the previous derivation, $\partial \loss/\partial z = f - y$, and by
the chain rule with $z = \vw\cdot\vx + b$,
\[
  \frac{\partial z}{\partial w_j} = x_j, \qquad \frac{\partial z}{\partial b} = 1 .
\]
Averaging over the training set,
\[
  \frac{\partial J}{\partial w_j}
  = \frac1m\sum_{i=1}^{m}\bigl(f^{(i)}-\yii{i}\bigr)\xjj{i}{j},
  \qquad
  \frac{\partial J}{\partial b}
  = \frac1m\sum_{i=1}^{m}\bigl(f^{(i)}-\yii{i}\bigr),
\]
or, in the vector form of \cref{sec:vector-form},
\[
  \grad_{\vw} J = \frac1m \Xmat^{\mathsf{T}}(\mathbf{f}-\vy),
  \qquad
  \frac{\partial J}{\partial b} = \frac1m \mathbf{1}^{\mathsf{T}}(\mathbf{f}-\vy),
  \qquad \mathbf{f} = \gsig(\Xmat\vw + b\mathbf{1}).
\]
\end{derivation}

Compare that with \cref{eq:vector-gradient}, the gradient for linear regression.
They are the same expression. Only the definition of $\mathbf{f}$ has changed,
from $\Xmat\vw + b\mathbf{1}$ to $\gsig(\Xmat\vw + b\mathbf{1})$.

\begin{algorithmbox}{Batch gradient descent for logistic regression}
\textbf{Input:} $\Xmat$ ($m\times n$), $\vy \in \{0,1\}^m$, learning rate
$\alpha$, iterations $N$.\\
\textbf{Output:} parameters $\vw, b$.

\begin{enumerate}\setlength{\itemsep}{1pt}
  \item Standardise every column of $\Xmat$, storing $\mu_j$ and $\sigma_j$
        (\cref{sec:scaling-practice}).
  \item $\vw \leftarrow \mathbf{0}$, \quad $b \leftarrow 0$.
  \item \textbf{for} $k = 1,\dots,N$:
  \begin{enumerate}\setlength{\itemsep}{1pt}
    \item $\mathbf{f} \leftarrow \gsig(\Xmat\vw + b\mathbf{1})$
          \hfill \emph{($O(mn)$)}
    \item $\mathbf{r} \leftarrow \mathbf{f} - \vy$
    \item $\vw \leftarrow \vw - \dfrac{\alpha}{m}\Xmat^{\mathsf{T}}\mathbf{r}$,
          \quad
          $b \leftarrow b - \dfrac{\alpha}{m}\sum_i r^{(i)}$
          \hfill \emph{(simultaneously)}
  \end{enumerate}
  \item \textbf{return} $\vw, b$.
\end{enumerate}

\textbf{Cost per iteration:} $O(mn)$, identical to
\cref{ch:ch04}. The only added work is one sigmoid per example per iteration.
\end{algorithmbox}

\begin{lstlisting}[style=pythonml,caption={Logistic regression. Line~6 is the
only line that differs from the linear-regression implementation of
\cref{ch:ch04}.}]
def sigmoid(z):
    return 1.0 / (1.0 + np.exp(-z))

def fit_logistic(X, y, alpha, iters):
    m, n = X.shape
    w, b = np.zeros(n), 0.0
    for _ in range(iters):
        f  = sigmoid(X @ w + b)            # (m,)  the only change
        r  = f - y                         # (m,)  same residual as before
        gw = (X.T @ r) / m                 # (n,)
        gb = r.mean()                      # scalar
        w, b = w - alpha * gw, b - alpha * gb     # simultaneous
    return w, b
\end{lstlisting}

\begin{pitfall}{Same update rule does not mean same algorithm}
Students routinely conclude from the identical formula that logistic regression
``is'' linear regression. It is not, and the differences are all examinable.

\begin{center}
\small
\renewcommand{\arraystretch}{1.25}
\begin{tabular}{@{}L{2.7cm}L{3.9cm}L{3.9cm}@{}}
\toprule
 & \textbf{linear regression} & \textbf{logistic regression}\\
\midrule
$f$        & $\vw\cdot\vx + b$ & $\gsig(\vw\cdot\vx + b)$\\
range of $f$ & all of $\R$ & $(0,1)$\\
cost $J$   & squared error & logistic loss\\
$J$ is     & quadratic and convex & convex, not quadratic\\
closed form & the normal equation & none exists\\
output means & the predicted value & $\Prob(y=1\mid\vx)$\\
\bottomrule
\end{tabular}
\end{center}

The formula is shared because both costs were derived from a probabilistic model
in the same exponential family, and in both the derivative of the loss with
respect to $z$ works out to $f - y$. That is a real and deep coincidence, not a
superficial one --- but it is a coincidence about the \emph{gradient}, and two
different functions can have gradients of the same shape.
\end{pitfall}

\begin{workedexample}{Two iterations by hand on T7}
Eight examples, one feature, $\alpha = 0.2$, starting from $w = b = 0$.

\textbf{Iteration 1.} At $w = b = 0$ every $z^{(i)} = 0$, so every
$f^{(i)} = \gsig(0) = 0.5$ and the cost is $J = \ln 2 = 0.693147$ exactly.

\begin{center}
\small
\renewcommand{\arraystretch}{1.15}
\begin{tabular}{@{}l*{8}{r}r@{}}
\toprule
$x$        & $1$ & $2$ & $3$ & $4$ & $5$ & $6$ & $7$ & $8$ & sum\\
$y$        & $0$ & $1$ & $0$ & $0$ & $0$ & $1$ & $1$ & $1$ & $4$\\
$f$        & $0.5$ & $0.5$ & $0.5$ & $0.5$ & $0.5$ & $0.5$ & $0.5$ & $0.5$ & \\
$f-y$      & $0.5$ & $-0.5$ & $0.5$ & $0.5$ & $0.5$ & $-0.5$ & $-0.5$ & $-0.5$ & $0$\\
$(f-y)x$   & $0.5$ & $-1.0$ & $1.5$ & $2.0$ & $2.5$ & $-3.0$ & $-3.5$ & $-4.0$ & $-5.0$\\
\bottomrule
\end{tabular}
\end{center}

Four of the eight are positive, so the residuals cancel and
$\partial J/\partial b = 0/8 = 0$. The weighted sum gives
$\partial J/\partial w = -5.0/8 = -0.625$, and the update is
\[
  w := 0 - 0.2\,(-0.625) = 0.125, \qquad b := 0 - 0.2\,(0) = 0 .
\]
The cost falls to $J = 0.663502$.

\textbf{Iteration 2.} Now $z^{(i)} = 0.125\,\xii{i}$, so the predictions spread
out --- but all of them are above $\tfrac12$, since $b = 0$ and every $x$ is
positive:

\begin{center}
\footnotesize
\setlength{\tabcolsep}{3pt}
\renewcommand{\arraystretch}{1.15}
\begin{tabular}{@{}l*{8}{r}r@{}}
\toprule
$x$      & $1$ & $2$ & $3$ & $4$ & $5$ & $6$ & $7$ & $8$ & sum\\
$f$      & $0.5312$ & $0.5622$ & $0.5927$ & $0.6225$ & $0.6514$ & $0.6792$
         & $0.7058$ & $0.7311$ & \\
$f-y$    & $0.5312$ & $-0.4378$ & $0.5927$ & $0.6225$ & $0.6514$ & $-0.3208$
         & $-0.2942$ & $-0.2689$ & $1.0759$\\
$(f-y)x$ & $0.5312$ & $-0.8756$ & $1.7780$ & $2.4898$ & $3.2568$ & $-1.9249$
         & $-2.0595$ & $-2.1515$ & $1.0442$\\
\bottomrule
\end{tabular}
\end{center}

So $\partial J/\partial b = 1.0759/8 = 0.134486$ and
$\partial J/\partial w = 1.0442/8 = 0.130526$, and
\[
  \begin{aligned}
    b &:= 0 - 0.2\,(0.134486) = -0.026897,\\
    w &:= 0.125 - 0.2\,(0.130526) = 0.098895 ,
  \end{aligned}
\]
with $J = 0.659160$. Both derivatives are now positive: the first step pushed
every prediction above $\tfrac12$, over-predicting on the four benign cases, and
the second step pulls the whole curve back down --- $b$ below zero, $w$ slightly
in from $0.125$. This is the elongated-valley behaviour of \cref{ch:ch05}, and
the next subsection measures it.
\end{workedexample}

\subsection{Everything from Chapter 5 still applies}

Run that same descent to convergence and it takes $1{,}461$ iterations, which is
a great many for eight points and two parameters. The reason is the one
\cref{sec:scaling-problem} identified, and the cure is the same:

\begin{center}
\small
\renewcommand{\arraystretch}{1.2}
\begin{tabular}{@{}lrrrrr@{}}
\toprule
features & $\lambda_{\min}$ & $\lambda_{\max}$ & $\kappa$ & $\alpha^\star$
         & iterations\\
\midrule
raw           & $0.04990$ & $6.57510$ & $131.8$ & $0.3042$ & $1{,}461$\\
standardised  & $0.25000$ & $0.25000$ & $1.0$   & $8.0000$ & $7$\\
\bottomrule
\end{tabular}
\end{center}

(Eigenvalues of the Hessian at the origin; each run uses $\alpha$ at $90\%$ of
its own stability limit and stops when the gradient norm falls below $10^{-6}$.)

The standardised condition number is exactly $1$, and for the same reason as in
\cref{sec:feature-scaling}: at the origin every $f^{(i)} = \tfrac12$, so every
weight $f(1-f)$ in the Hessian equals $\tfrac14$, and $H = \tfrac14 R$ where $R$
is the correlation matrix of the standardised design --- the identity, for a
single feature. The standardised optimum is $(b,w) = (0,\,1.361220)$, and
mapping back through $w = w_z/\sigma$, $b = b_z - w_z\mu/\sigma$ with
$\mu = 4.5$ and $\sigma = 2.291288$ recovers $w = 0.594085$ and $b = -2.673382$,
agreeing with the raw fit to six figures.

\begin{workedexample}{The fitted model on T7}
Converged parameters: $b = -2.673380$, $w = 0.594084$, $J = 0.528099$.
The decision boundary is $x^\star = -b/w = 4.5000$ cm, and each extra centimetre
multiplies the odds of malignancy by $e^{0.594084} = 1.8114$.

\begin{center}
\footnotesize
\setlength{\tabcolsep}{4pt}
\renewcommand{\arraystretch}{1.15}
\begin{tabular}{@{}l*{8}{r}@{}}
\toprule
$x$ (cm)         & $1$ & $2$ & $3$ & $4$ & $5$ & $6$ & $7$ & $8$\\
$y$              & $0$ & $1$ & $0$ & $0$ & $0$ & $1$ & $1$ & $1$\\
$f(x)$           & $0.1111$ & $0.1846$ & $0.2909$ & $0.4263$ & $0.5737$
                 & $0.7091$ & $0.8154$ & $0.8889$\\
prediction at $0.5$ & $0$ & $0$ & $0$ & $0$ & $1$ & $1$ & $1$ & $1$\\
\bottomrule
\end{tabular}
\end{center}

Six of eight correct. The two errors are the two examples that broke the
monotone pattern: the malignant $2$~cm tumour, which the model rates $0.1846$,
and the benign $5$~cm one, which it rates $0.5737$. Both are close calls in the
model's own terms, and \cref{sec:threshold} shows that the threshold repairs the
second for nothing and the first only at the price of three false alarms.
\end{workedexample}

\begin{pitfall}{Perfectly separable data has no maximum-likelihood fit}
Delete the awkward $2$~cm malignancy from T7 and the remaining seven points are
perfectly separated by $x = 4.5$. Now consider what maximum likelihood wants: it
can always reduce the cost further by scaling $\vw$ and $b$ up, pushing every
prediction closer to $0$ or $1$. The optimum is at infinity, the cost approaches
$0$ but never reaches it, and gradient descent produces coefficients that grow
for as long as you let it run:

\begin{center}
\small
\renewcommand{\arraystretch}{1.15}
\begin{tabular}{@{}rrrrr@{}}
\toprule
iterations & $b$ & $w$ & boundary & $J$\\
\midrule
$1{,}000$   & $-8.2117$  & $1.5282$ & $5.3733$ & $0.144837$\\
$10{,}000$  & $-20.3079$ & $3.7119$ & $5.4710$ & $0.042809$\\
$100{,}000$ & $-42.0652$ & $7.6646$ & $5.4882$ & $0.006149$\\
$200{,}000$ & $-49.5789$ & $9.0306$ & $5.4901$ & $0.003121$\\
\bottomrule
\end{tabular}
\end{center}

The boundary settles almost at once; the parameters never do.

You will meet this on small, clean, or high-dimensional datasets, and the symptom
is enormous coefficients with a suspiciously perfect training score. A solver
that reports convergence is reporting that it hit its iteration limit. The fix is the
penalty of \cref{ch:ch07}, which is why \tool{scikit-learn}'s
\code{LogisticRegression} applies an L2 penalty by default; \code{penalty=None}
turns it off and hands this failure straight back to you.
\end{pitfall}

\begin{researchspotlight}{Rosenblatt (1958) --- The Perceptron}
\textbf{The problem.} Build a machine that learns to classify patterns from
examples, without being programmed with the rule, and analyse when it will
succeed.

\textbf{The idea.} A single unit computes $z = \vw\cdot\vx + b$ and outputs a
class by thresholding $z$ at zero --- the decision boundary of
\cref{sec:decision-boundary}. Learning is an error-driven update: on a
misclassified example, $\vw \leftarrow \vw + \eta(y-\yhat)\vx$; on a correct one,
nothing.

\textbf{The result.} The perceptron convergence theorem: if the training data is
linearly separable, the procedure finds a separating hyperplane in finitely many
updates. If it is not separable --- as T7 is not --- the procedure never settles.

\textbf{Why it is here.} The perceptron and logistic regression have the same
hypothesis space and the same decision boundary, and their updates differ in one
respect: the perceptron thresholds first and then measures error, so
$y - \yhat \in \{-1,0,1\}$, while logistic regression measures error before
thresholding, so $y - f \in (-1,1)$. That single change is what buys a
differentiable cost, a probability, convergence on non-separable data, and the
possibility of stacking units into the networks of \cref{ch:ch08} --- where the
hard threshold has to go, precisely because it has no useful derivative.
\parencite{rosenblatt1958}
\end{researchspotlight}

\section{Choosing a Threshold}
\label{sec:threshold}

The threshold $0.5$ has been used throughout without an argument for it, and
there is none. It is the value at which the two classes are equally probable
under the model, which makes it the right choice only when the two kinds of
mistake cost the same. In screening they emphatically do not.

The two mistakes have names. A \term{false positive} is a negative example
predicted positive --- a healthy patient recalled for a biopsy. A
\term{false negative} is a positive example predicted negative --- a
malignancy sent home. Sweeping the threshold over the fitted T7 model trades one
against the other:

\begin{center}
\small
\renewcommand{\arraystretch}{1.2}
\begin{tabular}{@{}rrrrrrrr@{}}
\toprule
threshold & TP & FP & FN & TN & accuracy & precision & recall\\
\midrule
$0.10$ & $4$ & $4$ & $0$ & $0$ & $0.500$ & $0.500$ & $1.000$\\
$0.15$ & $4$ & $3$ & $0$ & $1$ & $0.625$ & $0.571$ & $1.000$\\
$0.30$ & $3$ & $2$ & $1$ & $2$ & $0.625$ & $0.600$ & $0.750$\\
$0.50$ & $3$ & $1$ & $1$ & $3$ & $0.750$ & $0.750$ & $0.750$\\
$0.60$ & $3$ & $0$ & $1$ & $4$ & $0.875$ & $1.000$ & $0.750$\\
$0.70$ & $3$ & $0$ & $1$ & $4$ & $0.875$ & $1.000$ & $0.750$\\
$0.85$ & $1$ & $0$ & $3$ & $4$ & $0.625$ & $1.000$ & $0.250$\\
$0.90$ & $0$ & $0$ & $4$ & $4$ & $0.500$ & --- & $0.000$\\
\bottomrule
\end{tabular}
\end{center}

Three things in that table are worth stopping on.

First, $0.5$ is not even the most accurate threshold. Raising it to $0.6$ turns
the benign $5$~cm tumour from a false positive into a correct rejection and
costs nothing, lifting accuracy from $0.750$ to $0.875$. The default was never
optimised for anything.

Second, no threshold gets everything right. The malignant $2$~cm tumour is rated
$0.1846$, below four benign cases; catching it requires a threshold under
$0.1846$, which by then admits three false positives. Thresholds reorder
decisions, they do not improve the model. Improving the model means better
features --- a texture score, as in T8 --- or a richer hypothesis space.

Third, the sensible way to choose is to price the two errors. If missing a
malignancy is judged as bad as ten unnecessary biopsies, minimise
$10 \times \text{FN} + 1 \times \text{FP}$. On this model that quantity is
minimised, at a cost of $3$, by any threshold in $(0.1111, 0.1846]$ --- the gap
between the predicted probabilities of the $1$~cm and $2$~cm tumours --- so
$t = 0.15$ is as good a choice as any, catching every malignancy at the price of
three false alarms.

\begin{keyinsight}{The model and the decision rule are separate}
Fitting produces a probability. Turning it into an action requires a threshold,
and the threshold is not part of the learning problem: it encodes the relative
cost of the two errors, which is a matter for the people who bear them, and it
can be changed after deployment without retraining.

Keep the two apart. Report the probabilities, choose the threshold from the
costs, and never let a library's default of $0.5$ stand in for a decision nobody
made. \Cref{ch:ch10} gives the metrics for judging the probabilities themselves,
and \cref{ch:ch11} the ones for judging thresholded decisions when the classes
are heavily imbalanced.
\end{keyinsight}

\begin{ethicsnote}{Who bears the false negatives?}
A threshold is a distributive decision. Moving it from $0.5$ to $0.15$ on the
model above converts one missed cancer into three unnecessary biopsies --- and
those are different people. The cost ratio that justifies the move is a value
judgement about whose harm counts and how much, dressed as a hyperparameter.

Two things follow for practice. State the ratio explicitly, in the model's
documentation, alongside who chose it; a threshold recorded only as
\code{0.5} in source code is a policy nobody can audit. And check the error rates
\emph{separately for each group the model will be applied to}: a single threshold
applied to a model whose probabilities are better calibrated for one group than
another silently gives the groups different error rates.
\Cref{ch:ch11} takes this up properly, with the measurements that make it
concrete.
\end{ethicsnote}

\begin{labbox}{Lab 6 --- Logistic regression from scratch}
Implement \code{sigmoid(z)} without overflow --- \code{np.exp(-z)} overflows for
$z$ around $-750$, so branch on the sign --- and verify
$\gsig'(z) = \gsig(z)(1-\gsig(z))$ against a finite-difference estimate to six
decimal places. Implement the logistic cost and gradient, and check the gradient
the same way.

Reproduce the T7 numbers: $J = \ln 2$ at the origin, the two hand-worked
iterations of \cref{sec:gd-logistic}, the converged $(b,w) =
(-2.673380,\,0.594084)$, and the boundary at exactly $4.5$~cm. Then fit the same
data with squared error from $(b,w) = (4,-1)$ and confirm that after $1{,}000$
iterations the sign of $w$ is still wrong. Finally, reproduce the threshold table
of \cref{sec:threshold} and plot precision and recall against the threshold on
one pair of axes.
\end{labbox}

\section*{Chapter Summary}
\addcontentsline{toc}{section}{Chapter Summary}

\begin{itemize}
  \item Classification needs a hypothesis space bounded to $[0,1]$ and a cost
        that charges for being wrong rather than far. Least squares on a $0/1$
        target fails both: it predicts $2.94$ for a $25$~cm tumour, and one
        distant correctly-classified example moves its boundary from $4.50$ to
        $5.57$.
  \item The sigmoid $\gsig(z) = 1/(1+e^{-z})$ maps $\R$ onto $(0,1)$, passes
        through $(0,\tfrac12)$, satisfies $1-\gsig(z) = \gsig(-z)$, and has
        derivative $\gsig' = \gsig(1-\gsig)$, whose maximum is $\tfrac14$ at
        $z=0$.
  \item The model $\fvwb(\vx) = \gsig(\vw\cdot\vx+b)$ is read as
        $\Prob(y=1\mid\vx)$. Inverting the sigmoid shows it is a linear model of
        the log-odds, so $e^{w_j}$ is the factor by which one unit of $x_j$
        multiplies the odds --- $1.8114$ per centimetre on T7.
  \item The decision boundary at threshold $\tfrac12$ is $\vw\cdot\vx + b = 0$,
        a hyperplane, and a property of the parameters alone. Polynomial
        features bend it into circles and beyond, at the cost of
        \cref{ch:ch07}'s problem.
  \item Squared error on a sigmoid model is non-convex --- both Hessian
        eigenvalues are negative at $(b,w) = (4,-1)$ on T7 --- and its gradient
        vanishes exactly where the model is most wrong: $0.000998$ for a
        prediction of $0.001$ against a target of $1$.
  \item Maximum likelihood on the model's own assumption yields the logistic
        loss $\loss = -y\ln f - (1-y)\ln(1-f)$ with no further choices. It is
        unbounded above, which is what keeps the gradient alive in the tails.
        At the origin $J = \ln 2 = 0.693147$ for every dataset.
  \item The logistic cost is convex, because $\dd^2\loss/\dd z^2 = f(1-f) > 0$
        and $z$ is affine in the parameters, and strictly convex when the
        augmented design matrix has full rank. It is not quadratic, so there is
        no normal equation and iteration is compulsory.
  \item The gradient is $\frac1m\Xmat^{\mathsf{T}}(\mathbf{f}-\vy)$: the same
        expression as linear regression with a different $\mathbf{f}$. Same form,
        different model, different cost, different guarantees.
  \item Feature scaling matters as much as before. On T7, standardising takes
        $\kappa$ from $131.8$ to $1.0$ and convergence from $1{,}461$ iterations
        to $7$.
  \item The threshold is a separate decision from the fit, chosen from the
        relative cost of false positives and false negatives. On T7, $0.6$ is
        more accurate than $0.5$, and $0.15$ is right if one missed malignancy
        is worth ten false alarms.
\end{itemize}

\section*{Exercises}
\addcontentsline{toc}{section}{Exercises}

\exercisesection{Short Questions}

\sq Give two independent reasons why a straight line fitted by least squares is
unsuitable for a binary target.

\sq State the sigmoid function and its four defining properties.

\sq Show that $1 - \gsig(z) = \gsig(-z)$, and say what this implies about
swapping the labels of the two classes.

\sq What is the log-odds, and in what sense is logistic regression linear?

\sq Define the decision boundary of a logistic model, and explain why it is a
property of the parameters and not of the data.

\sq Why does the sigmoid disappear from the decision rule when the threshold is
$\tfrac12$? What happens to the rule at a threshold of $0.7$?

\sq Give two reasons squared error is the wrong cost for a sigmoid model, one
about the shape of the cost and one about its gradient.

\sq Write the logistic loss in its two-case form and in its one-line form, and
show they agree for $y = 0$ and $y = 1$.

\sq What is the value of the logistic cost at $\vw = \mathbf{0}$, $b = 0$, and
why does it not depend on the data?

\sq Name three respects in which logistic regression differs from linear
regression, given that their gradient formulas are identical.

\sq A fitted logistic model has $w_j = -0.7$ for a feature measured in years.
State the effect of one additional year on the odds, and on the probability.

\exercisesection{Analytical Problems}

\ap Derive $\gsig'(z) = \gsig(z)(1-\gsig(z))$ from the definition, and hence show
that $\gsig'$ attains its maximum $\tfrac14$ at $z=0$ and nowhere else.

\ap Starting from $\Prob(y\mid\vx) = f^{y}(1-f)^{1-y}$, derive the logistic cost
of \cref{sec:logistic-loss} in full, justifying the logarithm, the negation and
the division by $m$.

\ap Show that $\partial \loss/\partial z = f - y$ for both $y=0$ and $y=1$, and
hence derive $\partial J/\partial w_j$ and $\partial J/\partial b$.

\ap Prove that the logistic cost is convex by showing that its Hessian is
$\frac1m\sum_i f^{(i)}(1-f^{(i)})\tilde{\vx}^{(i)}(\tilde{\vx}^{(i)})^{\mathsf{T}}$
and that this matrix is positive semi-definite. State the condition under which
it is positive definite, and relate it to \cref{sec:singular}.

\ap On T7 with $\alpha = 0.2$ and $w = b = 0$, verify by hand that
$\partial J/\partial b = 0$ and $\partial J/\partial w = -0.625$, and carry out
the first update. Then explain, without computing it, why
$\partial J/\partial b$ must be positive at the next step.

\ap A logistic model has $\vw = (2,-1)$ and $b = -4$. Write the decision boundary
in the form $x_2 = a + c\,x_1$, state which side is predicted positive, and find
the probability the model assigns to $\vx = (3, 1)$.

\ap Show that if the training data is linearly separable then the logistic cost
has no minimiser: for any $(\vw,b)$ that separates it, $J(2\vw, 2b) < J(\vw,b)$.
What does this imply about a solver that reports convergence on such data?

\exercisesection{Design Problems}

\dpr A bank asks you to build a model that decides which loan applications to
approve. Say what your model should output, what the threshold should be chosen
from, who should choose it, and what you would record so that the decision can
be audited two years later.

\dpr You are given a screening model whose accuracy is $99.2\%$ on a population
in which $0.8\%$ of cases are positive. Explain what is wrong with that report,
propose the measurements you would ask for instead, and describe the threshold
analysis you would run before recommending deployment.

\dpr A colleague's logistic regression returns coefficients of magnitude $10^{3}$
and perfect training accuracy on $40$ examples with $60$ features. Diagnose the
problem from this chapter alone, list the three checks you would run in order,
and say what you expect each to show.

\dpr Design the feature set for a logistic model predicting whether a customer
will cancel a subscription this month. Propose eight features, say what each
encodes, identify which need scaling and why, and state which pairs you would
expect to be strongly correlated.

\exercisesection{Programming Exercises}

\pe Implement a numerically stable \code{sigmoid} and verify the derivative
identity against a central finite difference at $z = -40, -1, 0, 1, 40$.
Explain why the naive implementation fails at $z = -800$ and yours does not.

\pe Implement the logistic cost and gradient and check the gradient against a
finite-difference estimate for a random $(\vw,b)$ on T8. Then reproduce the
converged fit of \cref{sec:decision-boundary}: $b = -9.827806$,
$w_1 = 0.931932$, $w_2 = 1.362882$, and the two misclassified points.

\pe Reproduce the divergence experiment of \cref{sec:squared-error-fails}: run
gradient descent on both costs from $(b,w) = (4,-1)$ with $\alpha = 0.2$ and
plot both cost curves on logarithmic axes. Report the iteration at which each
first attains $w > 0$.

\pe Reproduce the conditioning table of \cref{sec:gd-logistic} for T7, raw and
standardised, and confirm that mapping the standardised parameters back through
$w = w_z/\sigma$, $b = b_z - w_z\mu/\sigma$ recovers the raw fit to six figures.

\pe Fit a logistic model to T8 using the five quadratic features $x_1$, $x_2$,
$x_1^2$, $x_2^2$ and $x_1x_2$, and plot the resulting boundary by evaluating the
model on a grid and drawing its $0.5$ contour. Compare its training accuracy
with the linear model's, and say which you would deploy and why.

\exercisesection{Self-Test}

\qz The sigmoid $\gsig(z) = 1/(1+e^{-z})$ has range: (a)~$[-1,1]$; (b)~$(0,1)$;
(c)~$[0,1]$; (d)~all of $\R$.

\qz The derivative of the sigmoid is: (a)~$\gsig(z)$; (b)~$\gsig(z)(1-\gsig(z))$;
(c)~$1-\gsig(z)$; (d)~$e^{-z}$.

\qz The output of a logistic model is interpreted as: (a)~the predicted class;
(b)~$\Prob(y=1\mid\vx)$ under the model's assumption; (c)~the distance to the
boundary; (d)~the residual.

\qz Squared error applied to a sigmoid model gives a cost that is: (a)~convex;
(b)~non-convex, with vanishing gradients in the tails; (c)~quadratic;
(d)~undefined.

\qz The logistic loss for $y=1$ is: (a)~$(f-1)^2$; (b)~$-\ln f$; (c)~$-\ln(1-f)$;
(d)~$|f-1|$.

\qz The logistic cost at $\vw = \mathbf{0}$, $b=0$ equals: (a)~$0$; (b)~$0.25$;
(c)~$\ln 2 \approx 0.693$; (d)~it depends on the data.

\qz Logistic regression has a closed-form solution analogous to the normal
equation: (a)~yes, always; (b)~yes, when $\Xmat^{\mathsf{T}}\Xmat$ is invertible;
(c)~no, because the cost is convex but not quadratic; (d)~no, because the cost
is non-convex.

\qz The decision boundary of a logistic model at threshold $\tfrac12$ is:
(a)~$\gsig(\vw\cdot\vx+b) = 0$; (b)~$\vw\cdot\vx + b = 0$; (c)~$\vw\cdot\vx = 1$;
(d)~determined by the training data alone.

\qz Perfectly separable training data causes the fitted coefficients to:
(a)~become zero; (b)~grow without bound; (c)~stay unchanged; (d)~become complex.

\qz Raising the decision threshold from $0.5$ to $0.7$: (a)~always improves
accuracy; (b)~can only raise precision and can only lower recall;
(c)~retrains the model; (d)~changes the fitted parameters.

\furtherreading

\textcite{james2021} §4.1--4.3 is the primary reading for this chapter and covers
it almost exactly: §4.2 gives the argument of \cref{sec:linreg-fails} against
using linear regression for a binary response, and §4.3 develops the logistic
model, the logit, maximum-likelihood estimation and the reading of coefficients
as odds ratios, with a worked credit-default example that repays close study.
Read §4.3.4 for the multiple-feature case alongside \cref{sec:decision-boundary}.
\textcite{bishop2006} §4.3 is the supplementary treatment, and is worth the extra
effort for two things this chapter states without proving in full. Its §4.2
derives the sigmoid as the exact posterior of a natural class of probability
models, which answers \emph{why} this squashing function and not another; and
§4.3.3 derives the Hessian and the iteratively reweighted least squares
algorithm --- Newton's method on the same cost --- which is how the fits quoted
in \cref{sec:gd-logistic} were actually computed. §4.3.2 is the logistic model
itself, including the severe over-fitting on linearly separable data that is
this chapter's last pitfall.

\bigskip
\noindent\emph{Next:} \Cref{ch:ch07} names the phenomenon circling since
\cref{ch:ch05}'s degree-$9$ polynomial and this chapter's separable data ---
overfitting --- and supplies its standard defence, regularization, for both
models built so far.
